% Add "answers" option to print the answers
\documentclass[addpoints]{exam}

\usepackage[utf8]{inputenc}
\usepackage{calc}
\usepackage[dvipsnames]{xcolor}
\usepackage{ifthen}
\usepackage{tcolorbox}

% \makebox[\textwidth]{Name:\enspace\hrulefill}
%  
% \vspace{5mm}
%\makebox[\textwidth]{Instructor’s name:\enspace\hrulefill}

\definecolor{MyLightGray}{gray}{0.60}

\newcommand{\truefalse}
{		\begin{checkboxes}
			 \choice True
			 \choice False. Explanation: 
		\end{checkboxes}
		\vspace{10pt}
}


% Check if 'answers' is an option
\newcommand{\ifanswers}[2]{%
  \ifthenelse{\boolean{exam@answers}}{#1}{#2}%
}



\newcommand{\answerspace}[1]
{
	\ifprintanswers
	\else
		\setlength{\mylength}{10pt}%
		% \vspace{#1\mylength}
		% Make the grey box
		\par\nobreak\noindent
		\makebox[\linewidth][l]{%
			% set inner padding to zero, so only the frame shows
			{\setlength{\fboxsep}{0pt}%
			% draw a color frame, white (or transparent) inside
			\fcolorbox{MyLightGray}{white}{%
			% the inner box has zero width content but fixed height
			\parbox[t][#1\mylength][c]{%
				% subtract twice the frame‐thickness from the full width
				\dimexpr\linewidth-2\fboxrule\relax
			}{\strut}%
			}%
			}%
		}%
		\par
		\vspace{3\mylength}

  	\fi
}
\newlength{\mylength}

\newcommand{\code}[1]{{\texttt{#1}}}

\newcommand{\answer}[1]{\begin{solution}{#1}\end{solution}}




%\usepackage{listings}
\usepackage{color}
\usepackage{lastpage}
\usepackage{verbatim}
\usepackage{fancyvrb}
\usepackage{graphicx}
\usepackage{float}
\usepackage{minted}

\newcommand{\includecode}[2][c]{\lstinputlisting[caption=, escapechar=, language=#1]{#2}<!---->}
\RecustomVerbatimCommand{\VerbatimInput}{VerbatimInput}{}

\renewcommand{\solutiontitle}{\noindent\textbf{Solution sketch: }}

\definecolor{SolutionColor}{rgb}{0.8,0.9,1}

\cfoot{Page \thepage\ of \pageref{LastPage}}

\rhead{Initials: \rule{20mm}{0.4pt}}



%%%%%%%%%%%%%%%%%%%%%%%%%%%%%%%%%%%%%%%%
\begin{document}
\fvset{numbers=left,numbersep=3pt}
\thispagestyle{empty}

\begin{center}
	{\LARGE CS 408 Spring 2026 -- Practice Midterm \\}

	\vspace{20mm}
\fbox{\fbox{\parbox{5.5in}{\centering
	\begin{itemize}
          % \item Write your \textbf{name} in the first page and your \textbf{initials in every other page}.
        \item The exam is closed notes, closed book.

        \item The exam should be completed individually: any
          communication with your fellow students is strictly
          prohibited.

	\item Please print or write legibly.

        \item If you have a question, raise your hand and the
          course staff will come to you.

        \item Provide complete and clear but concise
          answers. Any explanations should be thorough.

        \item Answer the questions in the spaces provided. Your answer
          does not need to take up all the space provided; if you run
          out of room, use the empty pages at the back of the exam to
          complete your answer.

        \item Coding questions should be answered in Java, unless
          requested otherwise.

        \item Unless otherwise noted, CFG nodes represent basic blocks.
          % or low-level pseudo-code.
          % \item Unless otherwise noted, assume that questions refer to C code, AT\&T assembly syntax and the x86 architecture.

        \item Please double check that your copy of the exam
          is complete; you should have \pageref{LastPage}
          pages.

        \item You have 75 minutes to complete the exam.
        \end{itemize}
	}}}
\end{center}

\vspace{13mm}
\begin{center}
	\makebox[0.6\textwidth]{
		\combinedgradetable[h][questions]
	}
\end{center}
%\begin{center}\combinedgradetable[h][questions]\end{center}
\vspace{28mm}


\makebox[\textwidth]{\textbf{Name (please print)}: \rule{100mm}{0.8pt} }
%\end{center}

\vspace{10mm}
%\begin{center}
\makebox[\textwidth]{\textbf{PUID}: \rule{40mm}{0.8pt} \hspace{35mm}}
%\end{center}




\begin{questions}

	\newpage
%	\question State whether the following statements are true of false. \textbf{Explain the false statements.}
%	\begin{parts}
%		\part[2]
%		\truefalse
%		\answer{True.}
%
%	\end{parts}




\question Consider the concept of faults, errors and failures.

\begin{parts}

	\part[3] What is the difference between an error state and a failure? Explain.
	\answerspace{9}
	\answer{
          An error is an incorrect internal state of the program, such as a wrong value for a variable.
	}

      \part[3] What is the relationship between an error state and a failure? Explain.
	\answerspace{10}
	\answer{
          A failure is an external, incorrect behavior with respect to the expected behavior (e.g., crash, wrong output).
          An error is a necessary precondition for a failure.
          However, not all errors can cause a failure.
	}

\end{parts}

\question Consider the following three approaches to test a database management engine:
\begin{itemize}
	\item \textbf{Approach 1:} Run the target program on N
          randomly chosen inputs
	\item \textbf{Approach 2:} Run the target program on a
          representative set of N inputs
	\item \textbf{Approach 3:} Run the target program on all inputs
\end{itemize}

You can assume N is in the order of tens of thousands.

\begin{parts}

	\part[3] Explain the pros and cons of the approach 1 (N random inputs) and approach 3 (all inputs).
	\answerspace{11}
	\answer{

		Approach 1:

		\textbf{Pros:} Simple and efficient; can find some failures.

		\textbf{Cons:} Provides weak guarantees and likely misses corner cases.

		Approach 3:

		\textbf{Pros:} Comprehensive; will find every  fault that exists

		\textbf{Cons:} Unlikely to complete within the heat death of the universe.

	}

	\newpage
	\part[3] Describe how a a tester might use input space
          partitioning to identify a ``representative set of inputs''?
	\answerspace{9}
	\answer{
		A tester partitions the inputs into equivalence
                classes and selects inputs from each class.

		Inputs in the same class are expected to behave
                similarly, while inputs in different classes are
                expected to behave differently.
              }

	\part[3] What is the advantage of the approach 3 over the other two approaches (approach 1 and 2)?
	\answerspace{9}
	\answer{
		The main advantage is completeness.
		Approach 3 is an exhaustive test that can prove the absence of bugs.
		Approaches 1 and 2 are sampling techniques that can only prove the presence of bugs.
	}

\end{parts}


% \question Consider the RIP model, as discussed in class.

% \begin{parts}

% 	\part[5] What are the components of the RIP model? What is the relevance of the RIP model and each of its components for testing? Explain.
% 	\answerspace{12}
% 	\answer{
% 		RIP model defines three conditions that must be satisfied for a failure to be observed during testing, specifically:
% 		\begin{itemize}
% 			\item \textbf{Reachability:} The location or locations in the program that contain the fault must be reached.
% 			\item \textbf{Infection:} After executing the location, the state of the program must be incorrect.
% 			\item \textbf{Propagation:} The infected state must propagate to cause some output of the program to be incorrect.
% 		\end{itemize}
% 	}


% 	\part[3] What part is not covered by the RIP model that is also important for testing?
% 	\answerspace{12}
% 	\answer{
% 		Revealability. The tester must observe part of the incorrect output.
% 	}


% \end{parts}

\newpage

\question Consider the following function:
  \begin{minted}{Java}
public static int binarySearch(int[] arr, int x) {
        int l = 0;
        int r = arr.length - 1;
        while (l <= r) {
            int mid = l + (r - l) / 2;
            if (arr[mid] == x) {
                return mid;
            }
            if (arr[mid] < x) {
                l = mid + 1;
            } else {
                r = mid - 1;
            }
        }
        return -1; // element not found
    }

 \end{minted}

\begin{parts}

\part[4] Draw the CFG of this program. Label each node with the basic
  block it represents. Use a range of line numbers when necessary,
  e.g., use L2-7 as the label for a basic block starting at line 2
  and ending at line 7. If a line corresponds to more than one node,
  add a sequential letter suffix starting on ’a’ to distinguish
  them. For instance, if line 7 corresponds to two nodes CFG, you
  \textbf{must} name the first 7a and the second 7b.

  \answerspace{22}
  \answer{
    % \includegraphics[width=7cm]{midterm-q4-solution-cfg1.pdf}
  }


\part[3] What is a basic block?
  \answerspace{9}
  \answer{
    A basic block is a sequence of consecutive statements that has one entry point and one exit point, with no jump to/from the middle.
  }


\part[2] Is there any advantage in using basic blocks as opposed to statements in CFGs?
  \answerspace{9}
  \answer{
    Yes. It creates a much more compact graph, which is more efficient for control-flow and data-flow analyses.
  }

\part[2] Is there any disadvantage in using basic blocks as opposed to
  single statements in CFGs?
  \answerspace{9}
  \answer{
    For structural coverage analysis there is generally no disadvantage. (Other answers are also valid if well justified.)
    % Yes. It is impossible to analyze the program state between statements that are grouped into the same basic block.
  }

\end{parts}


\newpage

\question Consider the following Java function:
  \begin{minted}{Java}
public static int indexEven(int[] x, int n) {
 // if x is null, return -1
 // otherwise return the index of the n-th even number in array x
 // *or* return -1 if there is no n-th even number in array x
        if (x == null) {
            return -1;
        }
        int count = 0;
        for (int i = 0; i < x.length; i++) {
            if ((x[i] & 1) == 0) {
                count++;
                if (count = n) {
                   return i;
                 }
            }
        }
        return -1;
}
  \end{minted}

\begin{parts}
	\part[3] State where is the program fault and explain how you would fix the program. Consider only the smallest and simplest fix.
	\answerspace{10}
    \answer{
        % rubric:
        % location: 1
        % explain: 1
        % fix: 1
        %% Comment says the null pointer case is handled in exception.
        %One of the faults occurs in line 8, where the function does not check the a is null or not before the dereference. \\
        %Fix this fault by adding null checking before the for loop. \\
        Fault occurs in line 8, where the potentail out-of-bound access via the wrong condition in the for loop.
        Fix this fault by modifying \texttt{"i <= a.length"} to \texttt{"i < a.length"}.
    }

      \part[3] If possible, provide an example of a test input that
        does not cause the fault to execute.  If this is not possible
        explain why.

        \answerspace{10}
	\answer{
        % rubric:
        % Case 1 - have an example:
        %   - have: 1
        %   - input: 2
        % Case 2 - don't have an example:
        %   - don't have: 1
        %   - explain: 2
        It is impossible to avoid the fault in line 8,
        since there is no selection, jump, or iteration statements to let the program not to execute line 8.
    }

\newpage

  \part[3] If possible, provide an example of a test input that causes the fault to execute but no
    error to occur. If this is not possible explain why.
\answerspace{10}
    \answer{
        % rubric:
        % Case 1 - have an example:
        %   - have: 1
        %   - input satisfy requirement: 2
        % Case 2 - don't have an example:
        %   - don't have: 1
        %   - explain: 2
		% NOTE: Stating just the out-of-bounds or extra iteration is enough for full points on the justification.
        If \texttt{a = null}, the program will throw a \texttt{NullPointerException}, which is expected behavior. This will execute the fault but no error to occur.
        %There is no possible way to have a test input that causes the fault to execute but no error to occur.
        %The program will throw an out-of-bounds access exception or the program will always execute an extra iteration if the language runtime does not detect it. (Either justification is acceptable).
    }


  \part[3] If possible, provide an example of a test input that causes
    the fault to execute, the error to occur, but no failure to
    occur. If this is not possible explain why.

	\answerspace{10}
    \answer{
        % rubric:
        % Case 1 - have an example:
        %   - have: 1
        %   - input satisfy requirement: 2
        % Case 2 - don't have an example:
        %   - don't have: 1
        %   - explain: 2
        % In this case we have:
        % - provide invalid input: -2
        % - Not possible w/o oob access description: -2
        If the language does not detect the out-of-bounds access exception: \texttt{a = [ 1 ]},
        where the memory locations next to \texttt{a} is also zero.
		If the language detects the out-of-bounds access exception: not possible. (Only one explanation is required.)
    }


  \part[3] If possible, provide an example of a test input that causes
    the fault to execute, the failure to occur, but no error to
    occur. If this is not possible explain why.

    \answerspace{10}
    \answer{
        % rubric:
        % Case 1 - have an example:
        %   - have: 1
        %   - input satisfy requirement: 2
        % Case 2 - don't have an example:
        %   - don't have: 1
        %   - explain: 2
        % In this case we have:
        % - provide invalid input: -2
        % - Not possible w/o oob access description: -2
        If the language does not detect the out-of-bounds access exception: \texttt{a = [ 1 ]},
        where the memory locations next to \texttt{a} is also zero.
		If the language detects the out-of-bounds access exception: not possible. (Only one explanation is required.)
    }

	\newpage

	\part[4] Describe and explain the condition that an input must
          satisfy for the error to occur.
	\answerspace{10}
	\answer{
        % rebric:
        % input condition: 1
        % what/where error: 1
        % explain: 2
 		Any input will trigger an error due to the additional final iteration (i.e., wrong PC value). This applies regardless of whether the runtime implements memory safety or not.
        % The input condition is that \texttt{a} is allocated with a contiguous block of memory, including one extra element that holds a negative value in \texttt{int} type. \\
        % This will trigger the fault in line 8. \\
        % In the final iteration of the for loop, where \texttt{"i == a.length"}, the condition expression, \texttt{"i  <= a.length"}, is still true, resulting in out-of-bound memory access.
        % This behavior is incorrect.
        % However, \texttt{count} is still correct due to the extra element is a negative value. \\
        % If the program does not throw an out-of-bounds access exception, the expected value of \texttt{count} will still be returned by the function.
        % Thus, this input will trigger the error in the program.
    }

	\part[4] Describe and explain the condition that an input must
          satisfy for the failure to occur.

          \answerspace{10}
	\answer{
        % rebric:
        % input condition: 1
        % what/where failure: 1
        % explain: 2
		Two solutions accepted:
		\begin{enumerate}

		\item  If assuming memory-safety: Any input will trigger an out-of-bounds exception on the last iteration.

		\item If assuming no memory safety: Any input where the memory at a[a.length] is accessible and contains a value v that is greater than 0 and is odd. This will cause the variable ``count'' to increase beyond the expected result, which is returned.
		\end{enumerate}
        % The input condition is that \texttt{a} is allocated with a contiguous block of memory, including one extra element that holds a positive odd value in \texttt{int} type. \\
        % This will trigger the fault in line 8. \\
        % In the final iteration of the for loop, where \texttt{"i == a.length"}, the condition expression, \texttt{"i  <= a.length"}, is still true, resulting in \texttt{count} be incremented an addtional time.
        % This behavior is unintended, as \texttt{count} should not one extra time due to the out-of-bound access. \\
        % If the program does not throw an out-of-bounds access exception, the unexpected value of \texttt{count} will finally be returned by the function, which is an incorrect observed behavior.
        % Thus, this input will trigger the failure in the program. \\
        % If the program throw the out-of-bounds access, then any size of \texttt{a} will have failure. (Only one explanation is required.)
    }

\end{parts}

\newpage

\question Write below a single valid Java program with fewer than 12
  lines and 12 statements, that can not satisfy two different
  categories of coverage criteria that were discussed in class.

  \begin{figure}[ht]
    % \includegraphics[width=7cm]{midterm-cfg1.pdf}
  \end{figure}

\begin{parts}
\part[3] Identify the first coverage criterion that you selected and
  thoroughly explain why it is impossible for any test suite to
  completely satisfy that criteron for your program.

  \answerspace{7}
  \answer{
		Examples: [1, 3, 4, 5] and [1, 2, 4, 5].
		The test paths start from node 1 and end at node 5. There can be many possibilities of test paths in the given cfg, out of which, one path can be [1, 3, 4, 5] and another can be [1, 2 , 4, 5].
	}

      \part[3] Identify the second coverage criterion that you
        selected and thoroughly explain why it is impossible for any
        test suite to completely satisfy that criteron for your
        program.

\answerspace{7}
	\answer{
		[2],
		[2, 4],
		[2, 4, 5],
		[2, 4, 1]
		\\
		Any three of the above four simple paths are acceptable answers.
              }
            \end{parts}

\newpage


\newpage
\question Consider the following control-flow graph:
	\begin{figure}[ht]
          % \includegraphics[width=7cm]{midterm-cfg1.pdf}
	\end{figure}

\vspace{0.2cm}

\begin{parts}
	\part[3] If possible, list 2 test paths. If this is not possible, explain why.
	\answerspace{7}
	\answer{
		Examples: [1, 3, 4, 5] and [1, 2, 4, 5].
		The test paths start from node 1 and end at node 5. There can be many possibilities of test paths in the given cfg, out of which, one path can be [1, 3, 4, 5] and another can be [1, 2 , 4, 5].
	}

	\part[3] If possible, list 3 simple paths that (a) are not prime paths and (b) start on node 2. If this is not possible, explain why.
	\answerspace{7}
	\answer{
		[2],
		[2, 4],
		[2, 4, 5],
		[2, 4, 1]
		\\
		Any three of the above four simple paths are acceptable answers.
	}

	\part[3] If possible, list 3 prime paths that start on node 4. If this is not possible, explain why.
	\answerspace{7}
	\answer{
		There are only be 2 prime paths that start on 4 (which are [4, 1, 2, 4] and [4, 1, 3, 4]);  in total there are only 6 simple paths starting on 4, out of which only these two are not the proper subsets of other simple paths.
		}


	\newpage
      \part[10] Write an example of a syntactically valid Java function
        that generates the control-flow graph above. If this is not
        possible, explain why.


        \answerspace{23} \answer{ This is the closest possible
          solution- \VerbatimInput{midterm-q7_d_answer.txt} }

      \part[3] Provide test inputs that would generate the test paths
        you identified in 6a.


\end{parts}


\newpage
\question Consider the following control-flow graph annotated with
  def-use information. Unless otherwise noted, assume $def(l) =
  \emptyset$ and $use(l) = \emptyset$.
  \begin{figure}[ht]
    \begin{center}
      \includegraphics[width=.5\linewidth]{sample-midterm-cfg}
    \end{center}
\end{figure}


\begin{parts}
    \part[3] List the test requirement set of all-defs coverage with respect to \mintinline{Java}{x}.
    \answerspace{12}
    \answer{
		TR for all-defs coverage:
		TR = \{([0, 1, 2] or [0, 1, 4] or [0, 1, 2, 4])\}.
		% Accept answers with just
		%  TR = \{[0, 1, 2]\} OR
		%  TR = \{[0, 1, 4]\} OR
		%  TR = \{[0, 1, 2, 4]\}
	}

    \part[3] List the test requirement set of all-uses coverage with respect to \mintinline{Java}{x}.
    \answerspace{12}
    \answer{
		TR for all-uses coverage:
		TR = \{[0, 1, 2], ([0, 1, 4] or [0, 1, 2, 4])\}
		% Accept answer with just:
		%  - TR = {[0, 1, 2], [0, 1, 4]}    OR
        %  - TR = {[0, 1, 2], [0, 1, 2, 4]}

	}

    \part[4] List the test requirement set of all-du-paths coverage with respect to \mintinline{Java}{x}.
    \answerspace{12}
    \answer{
		TR for all-du-paths coverage:
		TR = \{[0, 1, 2], [0, 1, 4], [0, 1, 2, 4]\}
	}

\end{parts}

\newpage
\question[4]
  Provide a simple example of a program with a data race, where the
  lockset algorithm can help identify the fault. Provide and explain
  the results of the lockset algorithm on this program.

  \answerspace{20}
\answer{Two threads that write to a shared variable x without any locks. First access will have an empty candidate set (no locks held by the thread), leading to an alert.}

\question[4]
  Provide a simple example of a program without a problematic data
  race, where the lockset algorithm will nevertheless flag the data
  race for the user. Provide and explain the results of the lockset
  algorithm on this program, and argue why this is a false positive.

  \answerspace{20} \answer{Two threads that write to a shared variable
    x without any locks. First access will have an empty candidate set
    (no locks held by the thread), leading to an alert.}

\question[5] Explain how the feedback-based fuzzers discussed in class work.
\answerspace{20}
\answer{They use lightweight instrumentation to get the coverage (feedback) of each executed test. They iteratively generate new tests by mutating an initial or existing corpus test. If a test increases overall testing coverage, over the coverage of the corpus, it is added to the corpus.}


% \newpage
% \question[5] Describe two approaches to determine what are the properties of good bug reports. Explain how you would implement each.
% \answerspace{16}
% \answer{Examples: (a) conducting developer surveys and (b) studying the properties of bug reports in bug repositories.}


\newpage



\if 0 %%%%%%%%%%%%%%%%%%%%%%%%%%%%%%%%%%%%%%%%%
	\question[10]
	\answerspace{30}
	\answer{}

	\question
	\begin{parts}
		\part[6]
		\answerspace{20}
		\answer{}

		\part[2]
		\answerspace{10}
		\answer{}

	\end{parts}
\fi %%%%%%%%%%%%%%%%%%%%%%%%%%%%%%%%%%%%%%%%%

	\newpage
	\begin{center}
		[EMPTY PAGE 1]
	\end{center}

	\newpage
	\begin{center}
		[EMPTY PAGE 2]
	\end{center}

\end{questions}





\end{document}
